\documentclass{article}
\usepackage[margin=1in]{geometry}
\begin{document}
\title{JAL-SENSE Development Log}
\date{14 September 2026}
\maketitle
\section*{Version 1.0 -- MVP foundation}
\begin{itemize}
  \item Created a native Android Studio project using Kotlin and Jetpack Compose.
  \item Implemented Home, Dashboard, Alerts, Devices, History, and Settings navigation.
  \item Added simulated live sensor values, quality/status labels, treatment decision visibility, and equipment status.
  \item Added active/history alert views with acknowledge and resolve actions stored locally for the app session.
  \item Added data limitations and terms content, including the laboratory-validation disclaimer.
  \item Added a placeholder data boundary for future WebSocket/SSE or REST integration.
\end{itemize}
\section*{Known MVP limits}
This build is a UI prototype: authentication, persistent storage, backend ingestion, browser notifications, and custom-domain deployment require a backend/service deployment phase.
\section*{Verification}
Project files were checked for placeholder implementation markers and Android Studio project structure. A local Android SDK/Gradle runtime was not available in the workspace, so APK compilation should be run from Android Studio after its normal Gradle sync.
\section*{15 September 2026 -- Arduino sketch repair}
Corrected an incomplete SX1278 transmitter sketch supplied by the project owner. The source had been truncated inside the version-warning \texttt{Serial.printf} call, did not close \texttt{sendPacket()}, did not close \texttt{setup()}, and lacked a \texttt{loop()} function. A complete corrected sketch was saved as \texttt{JAL\_SENSE\_SX1278\_TX.ino}.
\section*{15 September 2026 -- Direct ESP32 dashboard integration}
Replaced simulated dashboard readings with a direct local-Wi-Fi ESP32 JSON client. The Android app requests \texttt{http://10.222.230.15/api/v1/measurements} every three seconds, parses its measurement array, shows connection failure safely, and lets the operator change the endpoint in Settings.
\end{document}
